\begin{alist}
\item Είναι:
\[1 - \log 2 = \log 10 - \log 2 = \log\frac{10}{2} = \log 5.\]

\item Επειδή $x^2 + 1 > 0$ η εξίσωση ορίζεται για κάθε $x \in \mathbb{R}$. Με τη βοήθεια του ερωτήματος $\alpha$) έχουμε:
\begin{align*}
\log(x^2 + 1) &= 1 - \log 2 \Leftrightarrow \\
\log(x^2 + 1) &= \log 5 \Leftrightarrow \\
x^2 + 1 &= 5 \Leftrightarrow \\
x^2 &= 4 \Leftrightarrow x = -2\text{ ή } x = 2.
\end{align*}
Άρα η εξίσωση έχει ρίζες τους αριθμούς $-2$ και $2$.
\end{alist}
